\chapter{Unit 3: Biodiversity and Conservation Biology}

\section{Section A: Multiple Choice Practice Questions (60 Curated Problems)}

\noindent\textbf{Q1. What does biodiversity refer to?} \hfill [Tag: Definition of biodiversity | Easy]
\begin{enumerate}[label=(\alph*),topsep=2pt,itemsep=1pt,partopsep=0pt,parsep=0pt]
    \item {The variety of living organisms}
    \item {The study of weather patterns}
    \item {The distribution of mineral resources}
    \item {The movement of tectonic plates}
\end{enumerate}
\vspace{2mm}

\noindent\textbf{Q2. Which level of biodiversity refers to variations in genes within a species?} \hfill [Tag: Levels of biological diversity: genetic, species and ecosystem diversity | Easy]
\begin{enumerate}[label=(\alph*),topsep=2pt,itemsep=1pt,partopsep=0pt,parsep=0pt]
    \item {Landscape diversity}
    \item {Species diversity}
    \item {Genetic diversity}
    \item {Ecosystem diversity}
\end{enumerate}
\vspace{2mm}

\noindent\textbf{Q3. The variety of different species present in an area is called:} \hfill [Tag: Levels of biological diversity: genetic, species and ecosystem diversity | Easy]
\begin{enumerate}[label=(\alph*),topsep=2pt,itemsep=1pt,partopsep=0pt,parsep=0pt]
    \item {Ecosystem diversity}
    \item {Species diversity}
    \item {Genetic diversity}
    \item {Cultural diversity}
\end{enumerate}
\vspace{2mm}

\noindent\textbf{Q4. Forests, grasslands, deserts, and wetlands found in a region demonstrate:} \hfill [Tag: Levels of biological diversity: genetic, species and ecosystem diversity | Easy]
\begin{enumerate}[label=(\alph*),topsep=2pt,itemsep=1pt,partopsep=0pt,parsep=0pt]
    \item {Ecosystem diversity}
    \item {Population diversity}
    \item {Species diversity}
    \item {Genetic diversity}
\end{enumerate}
\vspace{2mm}

\noindent\textbf{Q5. Which of the following is an ecosystem service provided by bees?} \hfill [Tag: Importance of biodiversity: ecosystem services | Easy]
\begin{enumerate}[label=(\alph*),topsep=2pt,itemsep=1pt,partopsep=0pt,parsep=0pt]
    \item {Pollination of crops}
    \item {Production of plastics}
    \item {Formation of rocks}
    \item {Extraction of metals}
\end{enumerate}
\vspace{2mm}

\noindent\textbf{Q6. The use of plants to develop medicines is an example of biodiversity's:} \hfill [Tag: Ecological, economic, social, aesthetic and informational value | Easy]
\begin{enumerate}[label=(\alph*),topsep=2pt,itemsep=1pt,partopsep=0pt,parsep=0pt]
    \item {Climatic value}
    \item {Astronomical value}
    \item {Economic value}
    \item {Tectonic value}
\end{enumerate}
\vspace{2mm}

\noindent\textbf{Q7. Enjoying the beauty of flowers and wildlife represents which value of biodiversity?} \hfill [Tag: Ecological, economic, social, aesthetic and informational value | Easy]
\begin{enumerate}[label=(\alph*),topsep=2pt,itemsep=1pt,partopsep=0pt,parsep=0pt]
    \item {Economic value}
    \item {Informational value}
    \item {Ecological value}
    \item {Aesthetic value}
\end{enumerate}
\vspace{2mm}

\noindent\textbf{Q8. What is a biodiversity hotspot?} \hfill [Tag: Hot spots of biodiversity | Easy]
\begin{enumerate}[label=(\alph*),topsep=2pt,itemsep=1pt,partopsep=0pt,parsep=0pt]
    \item {A region used mainly for intensive agriculture}
    \item {A region receiving very little rainfall}
    \item {A region containing many active volcanoes}
    \item {A region rich in endemic species and under threat}
\end{enumerate}
\vspace{2mm}

\noindent\textbf{Q9. Which activity is a direct cause of habitat loss?} \hfill [Tag: Threats to biodiversity: habitat loss | Easy]
\begin{enumerate}[label=(\alph*),topsep=2pt,itemsep=1pt,partopsep=0pt,parsep=0pt]
    \item {Clearing forests for farms}
    \item {Restoring degraded wetlands}
    \item {Planting native tree species}
    \item {Creating wildlife corridors}
\end{enumerate}
\vspace{2mm}

\noindent\textbf{Q10. What is wildlife poaching?} \hfill [Tag: Poaching of wildlife | Easy]
\begin{enumerate}[label=(\alph*),topsep=2pt,itemsep=1pt,partopsep=0pt,parsep=0pt]
    \item {Licensed treatment of injured animals}
    \item {Illegal hunting or capture of wild animals}
    \item {Seasonal movement of wild animals}
    \item {Scientific observation of wild animals}
\end{enumerate}
\vspace{2mm}

\noindent\textbf{Q11. What is an invasive alien species?} \hfill [Tag: Biological invasions | Easy]
\begin{enumerate}[label=(\alph*),topsep=2pt,itemsep=1pt,partopsep=0pt,parsep=0pt]
    \item {A migratory species that returns yearly}
    \item {A domestic species raised by people}
    \item {A native species with a small population}
    \item {A non-native species that causes harm}
\end{enumerate}
\vspace{2mm}

\noindent\textbf{Q12. Which situation is an example of human-wildlife conflict?} \hfill [Tag: Human-wildlife conflicts | Easy]
\begin{enumerate}[label=(\alph*),topsep=2pt,itemsep=1pt,partopsep=0pt,parsep=0pt]
    \item {Deer grazing within a national park}
    \item {Elephants damaging crops near a forest}
    \item {Butterflies pollinating flowers in a garden}
    \item {Birds nesting inside a protected wetland}
\end{enumerate}
\vspace{2mm}

\noindent\textbf{Q13. When is a species considered extinct?} \hfill [Tag: Species of conservation concern: extinct species | Easy]
\begin{enumerate}[label=(\alph*),topsep=2pt,itemsep=1pt,partopsep=0pt,parsep=0pt]
    \item {When no living individual remains}
    \item {When it moves to another habitat}
    \item {When it lives in one small region}
    \item {When its population begins decreasing}
\end{enumerate}
\vspace{2mm}

\noindent\textbf{Q14. An endangered species is one that:} \hfill [Tag: Endangered species | Easy]
\begin{enumerate}[label=(\alph*),topsep=2pt,itemsep=1pt,partopsep=0pt,parsep=0pt]
    \item {Has recently entered a new habitat}
    \item {Faces a very high risk of extinction}
    \item {Occurs commonly across many regions}
    \item {Is maintained mainly as a farm animal}
\end{enumerate}
\vspace{2mm}

\noindent\textbf{Q15. A vulnerable species is best described as a species that:} \hfill [Tag: Vulnerable species | Easy]
\begin{enumerate}[label=(\alph*),topsep=2pt,itemsep=1pt,partopsep=0pt,parsep=0pt]
    \item {Faces a high risk of extinction}
    \item {Exists only under human care}
    \item {Occurs abundantly in every habitat}
    \item {Has already disappeared completely}
\end{enumerate}
\vspace{2mm}

\noindent\textbf{Q16. Which feature is commonly associated with a rare species?} \hfill [Tag: Rare species | Easy]
\begin{enumerate}[label=(\alph*),topsep=2pt,itemsep=1pt,partopsep=0pt,parsep=0pt]
    \item {A complete absence of living individuals}
    \item {A rapid increase in every ecosystem}
    \item {A small population or limited distribution}
    \item {A large population across many continents}
\end{enumerate}
\vspace{2mm}

\noindent\textbf{Q17. What does EDGE stand for in biodiversity conservation?} \hfill [Tag: EDGE species | Easy]
\begin{enumerate}[label=(\alph*),topsep=2pt,itemsep=1pt,partopsep=0pt,parsep=0pt]
    \item {Evolutionarily Distinct and Globally Endangered}
    \item {Evolutionarily Dominant and Globally Endemic}
    \item {Environmentally Dependent and Genetically Enhanced}
    \item {Ecologically Diverse and Geographically Established}
\end{enumerate}
\vspace{2mm}

\noindent\textbf{Q18. Which is an example of in-situ conservation?} \hfill [Tag: Conservation of biodiversity: in-situ and ex-situ conservation | Easy]
\begin{enumerate}[label=(\alph*),topsep=2pt,itemsep=1pt,partopsep=0pt,parsep=0pt]
    \item {Breeding tigers in a zoological park}
    \item {Maintaining tissues in a laboratory}
    \item {Storing plant seeds in a seed bank}
    \item {Protecting tigers in a national park}
\end{enumerate}
\vspace{2mm}

\noindent\textbf{Q19. Which is an example of ex-situ conservation?} \hfill [Tag: Conservation of biodiversity: in-situ and ex-situ conservation | Easy]
\begin{enumerate}[label=(\alph*),topsep=2pt,itemsep=1pt,partopsep=0pt,parsep=0pt]
    \item {Protecting trees in a biosphere reserve}
    \item {Preserving plants in a sacred grove}
    \item {Conserving animals in a wildlife sanctuary}
    \item {Storing seeds in a seed bank}
\end{enumerate}
\vspace{2mm}

\noindent\textbf{Q20. How many major biogeographic zones are commonly recognized in India?} \hfill [Tag: Biogeographic zones of India | Easy]
\begin{enumerate}[label=(\alph*),topsep=2pt,itemsep=1pt,partopsep=0pt,parsep=0pt]
    \item {Ten}
    \item {Fourteen}
    \item {Twelve}
    \item {Eight}
\end{enumerate}
\vspace{2mm}

\noindent\textbf{Q21. A wetland survey records the variety of genes within fish populations, the number of plant and animal species, and the range of aquatic habitats. Which concept includes all three observations?} \hfill [Tag: Definition of biodiversity | Medium]
\begin{enumerate}[label=(\alph*),topsep=2pt,itemsep=1pt,partopsep=0pt,parsep=0pt]
    \item {Population density}
    \item {Biological productivity}
    \item {Biodiversity}
    \item {Ecological succession}
\end{enumerate}
\vspace{2mm}

\noindent\textbf{Q22. Farmers replace several traditional rice varieties with one high-yield variety, while the number of crop species grown remains unchanged. Which level of biodiversity declines most directly?} \hfill [Tag: Levels of biological diversity: genetic, species and ecosystem diversity | Medium]
\begin{enumerate}[label=(\alph*),topsep=2pt,itemsep=1pt,partopsep=0pt,parsep=0pt]
    \item {Species diversity}
    \item {Ecosystem diversity}
    \item {Landscape diversity}
    \item {Genetic diversity}
\end{enumerate}
\vspace{2mm}

\noindent\textbf{Q23. Two reserves contain the same number of species. Reserve P has similar population sizes for all species, while Reserve Q is dominated by one species. Which conclusion is most appropriate?} \hfill [Tag: Levels of biological diversity: genetic, species and ecosystem diversity | Medium]
\begin{enumerate}[label=(\alph*),topsep=2pt,itemsep=1pt,partopsep=0pt,parsep=0pt]
    \item {Reserve Q has greater species richness}
    \item {Reserve P has greater species evenness}
    \item {Reserve P has fewer ecological niches}
    \item {Reserve Q has greater genetic diversity}
\end{enumerate}
\vspace{2mm}

\noindent\textbf{Q24. A town restores mangroves instead of constructing an additional concrete barrier. During storms, the mangroves reduce wave energy and coastal flooding. Which ecosystem service is illustrated?} \hfill [Tag: Importance of biodiversity: ecosystem services | Medium]
\begin{enumerate}[label=(\alph*),topsep=2pt,itemsep=1pt,partopsep=0pt,parsep=0pt]
    \item {Regulating service}
    \item {Cultural service}
    \item {Supporting service}
    \item {Provisioning service}
\end{enumerate}
\vspace{2mm}

\noindent\textbf{Q25. Researchers study a wild plant and identify a chemical that can be developed into a new medicine. Which value of biodiversity is demonstrated most directly?} \hfill [Tag: Ecological, economic, social, aesthetic and informational value | Medium]
\begin{enumerate}[label=(\alph*),topsep=2pt,itemsep=1pt,partopsep=0pt,parsep=0pt]
    \item {Aesthetic value}
    \item {Informational value}
    \item {Recreational value}
    \item {Social value}
\end{enumerate}
\vspace{2mm}

\noindent\textbf{Q26. The decline of insect pollinators causes lower fruit production in nearby farms. Which combination of biodiversity values is most directly affected?} \hfill [Tag: Ecological, economic, social, aesthetic and informational value | Medium]
\begin{enumerate}[label=(\alph*),topsep=2pt,itemsep=1pt,partopsep=0pt,parsep=0pt]
    \item {Social and recreational values}
    \item {Aesthetic and informational values}
    \item {Ecological and economic values}
    \item {Social and aesthetic values}
\end{enumerate}
\vspace{2mm}

\noindent\textbf{Q27. Region X has many endemic plant species but retains most of its original vegetation. Region Y has many endemic plant species and has lost most of its original vegetation. Which region better fits the biodiversity hotspot concept?} \hfill [Tag: Hot spots of biodiversity | Medium]
\begin{enumerate}[label=(\alph*),topsep=2pt,itemsep=1pt,partopsep=0pt,parsep=0pt]
    \item {Region Y, because endemism and threat coincide}
    \item {Region Y, because all species are endangered}
    \item {Region X, because species abundance is higher}
    \item {Region X, because habitat must remain intact}
\end{enumerate}
\vspace{2mm}

\noindent\textbf{Q28. A highway divides a continuous forest into several isolated patches. Even if the total forest area changes only slightly, why may wildlife populations decline?} \hfill [Tag: Threats to biodiversity: habitat loss | Medium]
\begin{enumerate}[label=(\alph*),topsep=2pt,itemsep=1pt,partopsep=0pt,parsep=0pt]
    \item {Fragmentation increases food in every patch}
    \item {Isolation prevents competition among species}
    \item {Isolation eliminates all edge habitats}
    \item {Fragmentation restricts movement and gene flow}
\end{enumerate}
\vspace{2mm}

\noindent\textbf{Q29. Poachers repeatedly remove adult male elephants with large tusks from a population. Which long-term outcome is most likely?} \hfill [Tag: Threats to biodiversity: Poaching of wildlife | Medium]
\begin{enumerate}[label=(\alph*),topsep=2pt,itemsep=1pt,partopsep=0pt,parsep=0pt]
    \item {Reduced competition among plant species}
    \item {Improved gene flow between all herds}
    \item {Increased migration and larger herd size}
    \item {Reduced breeding and altered tusk traits}
\end{enumerate}
\vspace{2mm}

\noindent\textbf{Q30. An introduced aquatic plant spreads rapidly because local herbivores do not feed on it. It blocks sunlight and reduces native submerged plants. What best explains its success?} \hfill [Tag: Threats to biodiversity: Biological invasions | Medium]
\begin{enumerate}[label=(\alph*),topsep=2pt,itemsep=1pt,partopsep=0pt,parsep=0pt]
    \item {Mutualism with every native species}
    \item {Dependence on frequent forest fires}
    \item {Competitive release from natural enemies}
    \item {Low reproduction in disturbed habitats}
\end{enumerate}
\vspace{2mm}

\noindent\textbf{Q31. Elephants frequently raid crops in villages beside a fragmented reserve. Which intervention is most likely to reduce conflict while maintaining elephant movement?} \hfill [Tag: Threats to biodiversity: Human-wildlife conflicts | Medium]
\begin{enumerate}[label=(\alph*),topsep=2pt,itemsep=1pt,partopsep=0pt,parsep=0pt]
    \item {Moving all farms into the protected forest}
    \item {Blocking every route with permanent concrete walls}
    \item {Restoring corridors and using early-warning systems}
    \item {Removing all elephants that approach villages}
\end{enumerate}
\vspace{2mm}

\noindent\textbf{Q32. Repeated surveys find no members of a species in its former habitat, but a small breeding population survives in a zoo. Which status best describes it?} \hfill [Tag: Species of conservation concern: extinct species | Medium]
\begin{enumerate}[label=(\alph*),topsep=2pt,itemsep=1pt,partopsep=0pt,parsep=0pt]
    \item {Critically endangered}
    \item {Regionally vulnerable}
    \item {Extinct in the wild}
    \item {Extinct}
\end{enumerate}
\vspace{2mm}

\noindent\textbf{Q33. A species has undergone a severe population decline, occupies a shrinking range, and faces a very high risk of extinction in the wild. Which category is most appropriate?} \hfill [Tag: Species of conservation concern: Endangered species | Medium]
\begin{enumerate}[label=(\alph*),topsep=2pt,itemsep=1pt,partopsep=0pt,parsep=0pt]
    \item {Endangered}
    \item {Least concern}
    \item {Vulnerable}
    \item {Rare}
\end{enumerate}
\vspace{2mm}

\noindent\textbf{Q34. A reptile population is declining because nesting beaches are being developed, but its extinction risk is lower than that associated with the endangered category. How should it be classified?} \hfill [Tag: Species of conservation concern: Vulnerable species | Medium]
\begin{enumerate}[label=(\alph*),topsep=2pt,itemsep=1pt,partopsep=0pt,parsep=0pt]
    \item {Extinct}
    \item {Domesticated}
    \item {Vulnerable}
    \item {Data deficient}
\end{enumerate}
\vspace{2mm}

\noindent\textbf{Q35. A plant occurs naturally in only a few isolated valleys, but surveys show no current rapid decline. Which description fits it best?} \hfill [Tag: Species of conservation concern: Rare species | Medium]
\begin{enumerate}[label=(\alph*),topsep=2pt,itemsep=1pt,partopsep=0pt,parsep=0pt]
    \item {Extinct species}
    \item {Invasive species}
    \item {Rare species}
    \item {Common species}
\end{enumerate}
\vspace{2mm}

\noindent\textbf{Q36. Conservation funding must prioritize one of four equally endangered species. Which species would receive the highest priority under the EDGE approach?} \hfill [Tag: Species of conservation concern: EDGE species | Medium]
\begin{enumerate}[label=(\alph*),topsep=2pt,itemsep=1pt,partopsep=0pt,parsep=0pt]
    \item {A species with the widest geographic range}
    \item {A species from a highly distinct lineage}
    \item {A species with low public recognition}
    \item {A species from a large group of close relatives}
\end{enumerate}
\vspace{2mm}

\noindent\textbf{Q37. A critically threatened tree is protected in its native forest, while its seeds are also stored in a seed bank. How should these actions be classified?} \hfill [Tag: Conservation of biodiversity: in-situ and ex-situ conservation | Medium]
\begin{enumerate}[label=(\alph*),topsep=2pt,itemsep=1pt,partopsep=0pt,parsep=0pt]
    \item {Both actions are ex-situ conservation}
    \item {Forest protection is ex-situ; seed storage is in-situ}
    \item {Both actions are in-situ conservation}
    \item {Forest protection is in-situ; seed storage is ex-situ}
\end{enumerate}
\vspace{2mm}

\noindent\textbf{Q38. A species has only 12 individuals remaining, all scattered in a habitat facing immediate destruction. Which initial strategy offers the most direct short-term protection?} \hfill [Tag: Conservation of biodiversity: in-situ and ex-situ conservation | Medium]
\begin{enumerate}[label=(\alph*),topsep=2pt,itemsep=1pt,partopsep=0pt,parsep=0pt]
    \item {Waiting for natural population recovery}
    \item {Captive breeding with managed genetic records}
    \item {Unregulated tourism within its remaining habitat}
    \item {Introducing predators to restore food webs}
\end{enumerate}
\vspace{2mm}

\noindent\textbf{Q39. A conservation planner compares the cold high-altitude ecosystem of Ladakh with the warm, humid forests of the Western Ghats. Why are they treated as different biogeographic zones?} \hfill [Tag: Biogeographic zones of India | Medium]
\begin{enumerate}[label=(\alph*),topsep=2pt,itemsep=1pt,partopsep=0pt,parsep=0pt]
    \item {They have distinct climate, terrain, and biota}
    \item {They differ only in administrative boundaries}
    \item {They contain exactly the same endemic species}
    \item {They experience identical ecological processes}
\end{enumerate}
\vspace{2mm}

\noindent\textbf{Q40. Which combination best explains why India is recognized as a mega diversity nation?} \hfill [Tag: India as a mega diversity nation | Medium]
\begin{enumerate}[label=(\alph*),topsep=2pt,itemsep=1pt,partopsep=0pt,parsep=0pt]
    \item {Long coastline, one biome, and few native species}
    \item {High population, extensive farming, and urban growth}
    \item {Many ecosystems, high richness, and substantial endemism}
    \item {Large area, uniform climate, and low endemism}
\end{enumerate}
\vspace{2mm}

\noindent\textbf{Q41. A survey finds that two forest fragments contain the same number of tree species, but they differ greatly in allele frequencies, species abundances, and ecological interactions. Which conclusion best follows from the definition of biodiversity?} \hfill [Tag: Definition of biodiversity | Hard]
\begin{enumerate}[label=(\alph*),topsep=2pt,itemsep=1pt,partopsep=0pt,parsep=0pt]
    \item {Biodiversity cannot be compared unless every organism in both fragments is counted}
    \item {Only the fragment with greater total biomass can be considered more biodiverse}
    \item {Their biodiversity may differ because biodiversity includes variation within species and ecological organization}
    \item {The fragments have identical biodiversity because their species richness is equal}
\end{enumerate}
\vspace{2mm}

\noindent\textbf{Q42. Four isolated populations of a crop's wild relative are replaced by one high-yielding hybrid population containing the same number of individuals. Species richness and land-cover type remain unchanged. Which biodiversity component is most directly reduced?} \hfill [Tag: Levels of biological diversity: genetic, species and ecosystem diversity | Hard]
\begin{enumerate}[label=(\alph*),topsep=2pt,itemsep=1pt,partopsep=0pt,parsep=0pt]
    \item {Ecosystem diversity, because hybridization creates a new habitat}
    \item {Functional diversity, because individual abundance remains constant}
    \item {Genetic diversity, because locally adapted alleles may be lost}
    \item {Species diversity, because population isolation defines separate species}
\end{enumerate}
\vspace{2mm}

\noindent\textbf{Q43. Landscape P has 30 species in each of two habitats, with 25 species shared between them. Landscape Q also has 30 species in each habitat, but only 5 species are shared. Assuming comparable sampling, which statement is correct?} \hfill [Tag: Levels of biological diversity: genetic, species and ecosystem diversity | Hard]
\begin{enumerate}[label=(\alph*),topsep=2pt,itemsep=1pt,partopsep=0pt,parsep=0pt]
    \item {Q has higher beta diversity because species turnover is greater}
    \item {Q has lower ecosystem diversity because overlap is reduced}
    \item {P has higher beta diversity because its habitats share more species}
    \item {P has higher gamma diversity because local richness is identical}
\end{enumerate}
\vspace{2mm}

\noindent\textbf{Q44. A mangrove restoration project reduces cyclone damage, supports juvenile fish, stores carbon, and attracts visitors. Which mapping to ecosystem-service categories is most defensible?} \hfill [Tag: Importance of biodiversity: ecosystem services | Hard]
\begin{enumerate}[label=(\alph*),topsep=2pt,itemsep=1pt,partopsep=0pt,parsep=0pt]
    \item {Cyclone protection---provisioning; fish nursery---cultural; carbon storage---supporting; recreation---regulating}
    \item {Cyclone protection---supporting; fish nursery---regulating; carbon storage---provisioning; recreation---cultural}
    \item {Cyclone protection---regulating; fish nursery---supporting; carbon storage---regulating; recreation---cultural}
    \item {Cyclone protection---cultural; fish nursery---provisioning; carbon storage---regulating; recreation---supporting}
\end{enumerate}
\vspace{2mm}

\noindent\textbf{Q45. A wild plant has no current market use, but its genome contains an unusual drought-response pathway that may guide future crop breeding. Which value of biodiversity is most directly illustrated?} \hfill [Tag: Ecological, economic, social, aesthetic and informational value | Hard]
\begin{enumerate}[label=(\alph*),topsep=2pt,itemsep=1pt,partopsep=0pt,parsep=0pt]
    \item {Direct economic value arising from an established commercial market}
    \item {Aesthetic value arising from visual appreciation}
    \item {Social value arising from community identity}
    \item {Informational value arising from biological knowledge}
\end{enumerate}
\vspace{2mm}

\noindent\textbf{Q46. A wetland is valued at \text{Rs.~}40 million for harvested products and \text{Rs.~}90 million for avoided flood damage. Residents also oppose its destruction because it is sacred, even though they do not harvest from the sacred area. Which valuation is correct?} \hfill [Tag: Ecological, economic, social, aesthetic and informational value | Hard]
\begin{enumerate}[label=(\alph*),topsep=2pt,itemsep=1pt,partopsep=0pt,parsep=0pt]
    \item {\text{Rs.~}40 million and \text{Rs.~}90 million are both informational values, while sacredness is an ecosystem provisioning service}
    \item {\text{Rs.~}40 million is indirect ecological value, \text{Rs.~}90 million is social value, and sacredness is direct economic value}
    \item {\text{Rs.~}40 million is direct economic value, \text{Rs.~}90 million is indirect ecological value, and sacredness is social value}
    \item {\text{Rs.~}40 million is aesthetic value, \text{Rs.~}90 million is direct economic value, and sacredness is informational value}
\end{enumerate}
\vspace{2mm}

\noindent\textbf{Q47. Region X contains 1,600 endemic vascular plant species but retains 35\% of its original vegetation. Region Y contains 1,400 endemic vascular plant species and retains 20\%. Using the standard biodiversity-hotspot thresholds, which result follows?} \hfill [Tag: Hot spots of biodiversity | Hard]
\begin{enumerate}[label=(\alph*),topsep=2pt,itemsep=1pt,partopsep=0pt,parsep=0pt]
    \item {Only Y qualifies because it exceeds the habitat-loss threshold}
    \item {Both qualify because each satisfies at least one threshold}
    \item {Only X qualifies because it exceeds the endemic-plant threshold}
    \item {Neither qualifies because each fails one of the two thresholds}
\end{enumerate}
\vspace{2mm}

\noindent\textbf{Q48. A government must prioritize one of four regions under the biodiversity-hotspot concept. Which region best matches the concept's combined emphasis?} \hfill [Tag: Hot spots of biodiversity | Hard]
\begin{enumerate}[label=(\alph*),topsep=2pt,itemsep=1pt,partopsep=0pt,parsep=0pt]
    \item {A highly endemic region that has suffered extensive habitat loss}
    \item {A large wilderness with low endemism but exceptionally high total biomass and stable ecological processes}
    \item {A species-rich region with little endemism and intact habitat}
    \item {A degraded region containing many widespread generalist species}
\end{enumerate}
\vspace{2mm}

\noindent\textbf{Q49. A continuous forest is divided by roads into small patches, but the total forest area initially declines by only 5\%. Why can biodiversity loss still be disproportionately large?} \hfill [Tag: Threats to biodiversity: habitat loss | Hard]
\begin{enumerate}[label=(\alph*),topsep=2pt,itemsep=1pt,partopsep=0pt,parsep=0pt]
    \item {Fragmentation increases edge effects and isolates small populations}
    \item {Fragmentation automatically converts native species into invasive species}
    \item {Patch formation increases core habitat even when the total area remains nearly unchanged}
    \item {Roads eliminate genetic drift by increasing movement among populations}
\end{enumerate}
\vspace{2mm}

\noindent\textbf{Q50. Poachers selectively remove the largest tusked males from an elephant population over several generations. Which combined outcome is most plausible?} \hfill [Tag: Threats to biodiversity: poaching of wildlife | Hard]
\begin{enumerate}[label=(\alph*),topsep=2pt,itemsep=1pt,partopsep=0pt,parsep=0pt]
    \item {A rapid increase in genetic diversity and stabilization of the adult sex ratio}
    \item {A shift toward smaller tusks and disruption of age and social structure}
    \item {No evolutionary effect because poaching changes abundance but cannot alter heritable trait frequencies}
    \item {An increase in carrying capacity caused by stronger selection for large tusks}
\end{enumerate}
\vspace{2mm}

\noindent\textbf{Q51. An introduced plant becomes abundant only after a native herbivore is experimentally excluded. In invaded plots, native plant richness then falls sharply. Which inference is best supported?} \hfill [Tag: Threats to biodiversity: biological invasions | Hard]
\begin{enumerate}[label=(\alph*),topsep=2pt,itemsep=1pt,partopsep=0pt,parsep=0pt]
    \item {Herbivore exclusion proves that climate, propagule pressure, disturbance, and resource availability cannot influence the invasion}
    \item {Native richness declines because all introduced species inevitably hybridize with every native species}
    \item {The introduced plant is harmless because its abundance depends on herbivore exclusion}
    \item {Enemy release may facilitate invasion, followed by competitive suppression of native plants}
\end{enumerate}
\vspace{2mm}

\noindent\textbf{Q52. Elephant crop-raiding rises after a traditional movement corridor is converted into farms. Which intervention most directly addresses the underlying landscape driver while limiting harm to both people and elephants?} \hfill [Tag: Threats to biodiversity: human-wildlife conflicts | Hard]
\begin{enumerate}[label=(\alph*),topsep=2pt,itemsep=1pt,partopsep=0pt,parsep=0pt]
    \item {Relocate all elephants without assessing habitat capacity at the release site}
    \item {Provide compensation only after every verified incident}
    \item {Restore corridor connectivity and combine it with locally managed crop-protection measures}
    \item {Build continuous impermeable barriers around every protected area, settlement, farm, river crossing, and seasonal range}
\end{enumerate}
\vspace{2mm}

\noindent\textbf{Q53. A freshwater fish has not been recorded for 40 years despite repeated surveys. Several deep, inaccessible tributaries within its former range remain unsurveyed. Which classification judgment is scientifically strongest?} \hfill [Tag: Species of conservation concern: extinct species | Hard]
\begin{enumerate}[label=(\alph*),topsep=2pt,itemsep=1pt,partopsep=0pt,parsep=0pt]
    \item {Classify it as extinct in the wild because no captive population is known}
    \item {Declare it extinct because four decades without a record is sufficient}
    \item {Classify it as least concern because unsurveyed habitat remains}
    \item {Avoid declaring it extinct until reasonable doubt about surviving individuals is removed}
\end{enumerate}
\vspace{2mm}

\noindent\textbf{Q54. Two species face extinction risk. Species A has undergone a rapid, severe population decline across most of its range. Species B has a stable, large population but modest habitat deterioration. Which statement best distinguishes an Endangered assessment from ordinary rarity?} \hfill [Tag: Species of conservation concern: Endangered species | Hard]
\begin{enumerate}[label=(\alph*),topsep=2pt,itemsep=1pt,partopsep=0pt,parsep=0pt]
    \item {Species A may qualify because Endangered status reflects very high extinction risk, not rarity alone}
    \item {Species B must qualify because every habitat change produces Endangered status}
    \item {Species A cannot qualify unless fewer than 50 individuals remain}
    \item {Both qualify automatically because conservation categories are based only on expert preference}
\end{enumerate}
\vspace{2mm}

\noindent\textbf{Q55. A species does not meet the criteria for Endangered but satisfies validated criteria indicating a high risk of extinction in the wild. Which category is most appropriate?} \hfill [Tag: Species of conservation concern: Vulnerable species | Hard]
\begin{enumerate}[label=(\alph*),topsep=2pt,itemsep=1pt,partopsep=0pt,parsep=0pt]
    \item {Vulnerable}
    \item {Rare}
    \item {Endangered}
    \item {Extinct in the Wild}
\end{enumerate}
\vspace{2mm}

\noindent\textbf{Q56. A naturally scarce plant occurs in small populations scattered across a broad area. Its abundance is stable, and no major threat is currently documented. Why should "rare" not automatically be treated as equivalent to "Endangered"?} \hfill [Tag: Species of conservation concern: Rare species | Hard]
\begin{enumerate}[label=(\alph*),topsep=2pt,itemsep=1pt,partopsep=0pt,parsep=0pt]
    \item {Rarity is determined solely by body size, while threat categories are determined solely by economic value}
    \item {Endangered species must be common but declining, whereas rare species must be stable}
    \item {Rarity describes scarcity, whereas Endangered denotes an assessed very high extinction risk}
    \item {Rare species are necessarily secure because broad distribution prevents local extinction}
\end{enumerate}
\vspace{2mm}

\noindent\textbf{Q57. Species P is the sole surviving member of an ancient lineage and is globally Endangered. Species Q is Critically Endangered but belongs to a recently diversified lineage with many close relatives. Why might P receive a high EDGE priority despite Q's greater immediate threat?} \hfill [Tag: Species of conservation concern: EDGE species | Hard]
\begin{enumerate}[label=(\alph*),topsep=2pt,itemsep=1pt,partopsep=0pt,parsep=0pt]
    \item {EDGE rankings ignore extinction risk when a species has no close relatives}
    \item {Q cannot be prioritized because recently evolved species have no ecological value}
    \item {P combines substantial extinction risk with unusually high evolutionary distinctiveness}
    \item {P is automatically more abundant because its lineage is ancient}
\end{enumerate}
\vspace{2mm}

\noindent\textbf{Q58. A tree species is reduced to 12 wild individuals, its seeds cannot survive conventional seed-bank drying, and its habitat is still being degraded. Which strategy is most defensible?} \hfill [Tag: Conservation of biodiversity: in-situ and ex-situ conservation | Hard]
\begin{enumerate}[label=(\alph*),topsep=2pt,itemsep=1pt,partopsep=0pt,parsep=0pt]
    \item {Transfer the species permanently to ornamental gardens and abandon protection of its ecological interactions and remaining habitat}
    \item {Use habitat protection alone because ex-situ methods always reduce survival}
    \item {Protect and restore the habitat while establishing living collections or cryopreserved tissues}
    \item {Use a conventional dry-seed bank alone because all plant seeds tolerate desiccation}
\end{enumerate}
\vspace{2mm}

\noindent\textbf{Q59. A conservation portfolio includes cold high-altitude plateaus of Ladakh, the alluvial plains of the Ganga, coral-bearing oceanic islands, and the rain-shadow plateau east of the Western Ghats. Which sequence of Indian biogeographic zones best matches these areas?} \hfill [Tag: Biogeographic zones of India | Hard]
\begin{enumerate}[label=(\alph*),topsep=2pt,itemsep=1pt,partopsep=0pt,parsep=0pt]
    \item {Himalaya, Semi-arid, Coasts, Western Ghats}
    \item {Trans-Himalaya, Gangetic Plain, Islands, Deccan Peninsula}
    \item {Himalaya, Gangetic Plain, Northeast, Deccan Peninsula}
    \item {Trans-Himalaya, Desert, Coasts, Gangetic Plain}
\end{enumerate}
\vspace{2mm}

\noindent\textbf{Q60. Which combination most strongly supports India's recognition as a megadiverse nation rather than merely a country with a large land area?} \hfill [Tag: India as a mega diversity nation | Hard]
\begin{enumerate}[label=(\alph*),topsep=2pt,itemsep=1pt,partopsep=0pt,parsep=0pt]
    \item {High species richness and endemism across sharply contrasting biogeographic zones}
    \item {Uniform climate and low habitat turnover across the entire subcontinent}
    \item {A large human population and extensive cultivation across several river basins}
    \item {The complete inclusion of every global biodiversity hotspot within India's political boundaries}
\end{enumerate}
\vspace{2mm}


\section{Section B: Comprehensive Subjective Problem with Detailed Analytical Rubric}

\begin{subjectiveq}{3}{1}
\textbf{Quantitative Biodiversity Indices and Conservation Prioritization:}\\An ecological quadrant census of an undisturbed primary moist evergreen forest patch in the Silent Valley (Western Ghats Hotspot) records five endemic arboreal species with the following observed counts of individuals:\begin{center}\begin{tabular}{llc}\toprule\textbf{Species Identifier} & \textbf{Scientific Designation} & \textbf{Individual Count ($n_i$)} \\\midrule Species A & \textit{Cullenia exarillata} & 64 \\Species B & \textit{Palaquium ellipticum} & 48 \\Species C & \textit{Mesua ferrea} & 28 \\Species D & \textit{Calophyllum polyanthum} & 16 \\Species E & \textit{Dipterocarpus bourdillonii} & 4 \\\bottomrule\end{tabular}\end{center}(a) Compute the total abundance $N$ and the relative abundance proportion $p_i = n_i / N$ for each species.\\(b) Calculate the Shannon-Wiener Diversity Index: $H' = -\sum_{i=1}^S p_i \ln(p_i)$.\\(c) Calculate the Maximum Theoretical Diversity $H'_{\text{max}} = \ln(S)$ and the Pielou's Species Evenness Index $J' = H' / H'_{\text{max}}$.\\(d) Compute Simpson's Dominance Index $D = \sum_{i=1}^S p_i^2$ and Simpson's Index of Diversity ($1 - D$).\\(e) \textit{Dipterocarpus bourdillonii} is categorized as Critically Endangered (CR) under IUCN criteria. Formulate a comparative analysis of in-situ (biosphere reserve core zoning) versus ex-situ (cryogenic seed banking, arboreta) preservation strategies for this species.
\end{subjectiveq}

\begin{solbox}
\textbf{Part (a): Abundance and Proportions:}\\Total individual count $N = 64 + 48 + 28 + 16 + 4 = 160$ across $S = 5$ species.\begin{align*}p_A &= \frac{64}{160} = 0.40, \quad p_B = \frac{48}{160} = 0.30, \quad p_C = \frac{28}{160} = 0.175 \\p_D &= \frac{16}{160} = 0.10, \quad p_E = \frac{4}{160} = 0.025.\end{align*}\textbf{Part (b): Shannon-Wiener Index $H'$:}\\Evaluating $p_i \ln(p_i)$ for each species:\begin{align*}p_A \ln(p_A) &= 0.40 \times \ln(0.40) = 0.40 \times (-0.91629) = -0.36652 \\p_B \ln(p_B) &= 0.30 \times \ln(0.30) = 0.30 \times (-1.20397) = -0.36119 \\p_C \ln(p_C) &= 0.175 \times \ln(0.175) = 0.175 \times (-1.74297) = -0.30502 \\p_D \ln(p_D) &= 0.10 \times \ln(0.10) = 0.10 \times (-2.30259) = -0.23026 \\p_E \ln(p_E) &= 0.025 \times \ln(0.025) = 0.025 \times (-3.68888) = -0.09222.\end{align*}Summing all components:\begin{equation*}\sum_{i=1}^5 p_i \ln(p_i) = -1.35521 \implies H' = -(-1.35521) = 1.3552\text{ nats}.\end{equation*}\textbf{Part (c): Maximum Diversity and Evenness $J'$:}\\\begin{align*}H'_{\text{max}} &= \ln(S) = \ln(5) = 1.60944 \\J' &= \frac{H'}{H'_{\text{max}}} = \frac{1.35521}{1.60944} = 0.8420\text{ (or } 84.20\%\text{ evenness)}.\end{align*}\textbf{Part (d): Simpson's Dominance and Diversity Indices:}\\\begin{align*}D = \sum p_i^2 &= (0.40)^2 + (0.30)^2 + (0.175)^2 + (0.10)^2 + (0.025)^2 \\&= 0.1600 + 0.0900 + 0.030625 + 0.0100 + 0.000625 = 0.29125 \\\text{Simpson's Index of Diversity} &= 1 - D = 1 - 0.29125 = 0.70875.\end{align*}\textbf{Part (e): Conservation Strategy Comparison:}\\\begin{itemize}[leftmargin=1.5em]  \item \textbf{In-Situ Protection:} Preservation of the microclimate, mycorrhizal soil fungal networks, and native pollinators (bats, bees) in the core zone of Nilgiri Biosphere Reserve. Prohibits road construction, linear corridor fragmentation, and cattle grazing.  \item \textbf{Ex-Situ Backup:} Collection of recalcitrant seeds for cryogenic tissue culture at the Jawaharlal Nehru Tropical Botanic Garden and Research Institute (TBGRI), establishing insurance populations to mitigate catastrophic wildfire risk.\end{itemize}
\end{solbox}

\begin{examinertip}
\textbf{Marking Scheme Breakdown (Total: 10 Marks):}\\
\textbullet\ \textbf{3 Marks:} Clear formulation of governing environmental identities, equations, and thermodynamic/constitutional foundations.\\
\textbullet\ \textbf{4 Marks:} Rigorous algebraic/numerical substitution, dimensional consistency, and unit handling.\\
\textbullet\ \textbf{3 Marks:} Critical evaluation of environmental implications, regulatory criteria (CPCB/UNDRR/IUCN), and policy mitigations.
\end{examinertip}
